\clearpage
\section*{Problem 2}
\addcontentsline{toc}{section}{Problem 2}
\begingroup
\hypersetup{colorlinks=true,
            linkcolor=blue!50!black,
            citecolor=blue!50!black,
            urlcolor=blue!50!black}
\newcommand{\qtwoF}{F}
\newcommand{\qtwooo}{\mathfrak{o}}
\newcommand{\qtwopp}{\mathfrak{p}}
\newcommand{\qtwoqq}{\mathfrak{q}}
\newcommand{\qtwoC}{\mathbb{C}}
\newcommand{\qtwoZ}{\mathbb{Z}}
\newcommand{\qtwoGL}{\mathrm{GL}}
\newcommand{\qtwoMn}{M_n}
\newcommand{\qtwoSe}{\mathcal{S}_{\mathrm e}}
\newcommand{\qtwoSw}{\mathcal{S}}
\newcommand{\qtwoW}{\mathcal{W}}
\newcommand{\qtwovol}{\operatorname{vol}}
\newcommand{\qtwodiag}{\operatorname{diag}}
\newcommand{\qtwotr}{\operatorname{trace}}
\newcommand{\qtwottilde}[1]{\widetilde{#1}}
\newcommand{\qtwolRS}{\ell_{\mathrm{RS}}}
\newtheorem{qtwotheorem}{Theorem}
\newtheorem{qtwolemma}{Lemma}
\newtheorem{qtwoproposition}{Proposition}
\newtheorem{qtwocorollary}{Corollary}
\title{Solution to Problem 2: A Universal Test Vector\\
for Rankin--Selberg Integrals}
\author{}
\date{}
\maketitle


\section*{Statement and Answer}

Let $\qtwoF$ be a non-archimedean local field with ring of integers $\qtwooo$, maximal
ideal $\qtwopp$, residue cardinality $q$, and a nontrivial additive character
$\psi:\qtwoF\to\qtwoC^\times$ of conductor $\qtwooo$. Write $G_r:=\qtwoGL_r(\qtwoF)$, let $N_r<G_r$
be the upper-triangular unipotent subgroup, and embed $G_n\hookrightarrow
G_{n+1}$ as the upper-left block. Set $K_r:=\qtwoGL_r(\qtwooo)$, with Haar measures
giving $K_r$ and $N_r\cap K_r$ volume one.

\medskip
\noindent\textbf{Theorem.} \emph{Let $\Pi$ be a generic irreducible admissible
representation of $G_{n+1}$, realized in its $\psi^{-1}$-Whittaker model
$\qtwoW(\Pi,\psi^{-1})$. Then there exists a single $W\in\qtwoW(\Pi,\psi^{-1})$ with the
following property. For every generic irreducible admissible $\pi$ of $G_n$ with
conductor ideal $\qtwoqq$, generator $Q$ of $\qtwoqq^{-1}$, and
$u_Q:=I_{n+1}+QE_{n,n+1}$, there exists $V\in\qtwoW(\pi,\psi)$ such that}
\[
\int_{N_n\backslash G_n} W(\qtwodiag(g,1)u_Q)\,V(g)\,|\det g|^{s-\frac12}\,dg
\]
\emph{is finite and nonzero for all $s\in\qtwoC$.}

\medskip
\noindent The answer is \textbf{YES}. The single $W$ may be taken to be an
explicit, $\pi$-independent Kirillov-model vector, and for each $\pi$
the role of $V$ is played by the \emph{normalized newvector}. The integral is in
fact \emph{constant in $s$} and equal to an explicit nonzero scalar; this is why
finiteness and nonvanishing hold for all $s$, with no analytic continuation or
pole-cancellation required.

\medskip
\noindent\textbf{A failure mode to avoid.} One cannot choose $V$ as a
Schwartz bump in the Kirillov model making the integrand constant on its
support: $V$ has a central character, so it cannot be compactly supported modulo
the center. (Already for $n=1$ the relevant integral is a normalized Gauss sum,
whose integrand is genuinely non-constant.) The correct argument runs through the
Godement--Jacquet functional equation, not through test-vector flexibility.

\section*{Choice of $W$ and $V$}

For $x\in\qtwoF$ and $y\in\qtwoF^\times$ write
\[
d_y:=\qtwodiag(1,\dots,1,y)\in G_n\hookrightarrow G_{n+1},\qquad
u_x:=I_{n+1}+xE_{n,n+1}\in N_{n+1}.
\]

\noindent\textbf{The vector $W$.} Define $W_0$ on $G_n$ by
\begin{equation}\label{eq:W0}
W_0(g):=\int_{N_n}\mathbf 1_{K_n}(xg)\,\psi(x)\,dx,
\end{equation}
a function on $G_n$, left $(N_n,\psi^{-1})$-equivariant, which extends to an element of
$\qtwoW(\Pi,\psi^{-1})$ on $G_{n+1}$ by Lemma~\ref{lem:extend} below. We take
$W:=u_Q W_0$ (equivalently we evaluate this extension at the shifted argument); the
extension, hence $W$, depends only on $\Pi$, not on $\pi$. Throughout we identify $G_n$
with its image in $G_{n+1}$ under $g\mapsto\qtwodiag(g,1)$, and we use the same symbol
$W_0$ for the function on $G_n$ in \eqref{eq:W0} and for its chosen extension to
$G_{n+1}$; this is justified by Lemma~\ref{lem:extend}, which guarantees that the
extension restricts to \eqref{eq:W0} on $\qtwodiag(G_n,1)$.

The function $W_0$ on $G_n$ is supported on $N_nK_n$ and satisfies $W_0(k)=1$ for
$k\in K_n$. Indeed $W_0(g)\neq0$ forces $xg\in K_n$ for some $x\in N_n$, i.e.\
$g\in N_nK_n$; for $g=nk$ with $n\in N_n$, $k\in K_n$ the left $\psi^{-1}$-equivariance
gives $W_0(nk)=\psi^{-1}(n)W_0(k)$, and $W_0(k)=\int_{N_n}\mathbf 1_{K_n}(xk)\psi(x)\,dx
=\int_{N_n\cap K_n}\psi(x)\,dx=1$ because $\psi$ is trivial on $N_n\cap K_n$ (conductor
$\qtwooo$) and $\qtwovol(N_n\cap K_n)=1$.

\medskip
We now justify the two facts about $W_0$ used below: that \eqref{eq:W0} is the
restriction to $\qtwodiag(G_n,1)$ of a genuine element of $\qtwoW(\Pi,\psi^{-1})$, and that
this extension satisfies the right-$u_1$ translation identity. Let
$P_{n+1}:=\{p\in G_{n+1}:\text{last row of }p\text{ is }(0,\dots,0,1)\}$ be the
mirabolic subgroup of $G_{n+1}$; it is generated by $\qtwodiag(G_n,1)$ and $N_{n+1}$, and
$\qtwodiag(G_n,1)$ normalizes the abelian unipotent radical
$U:=\{I_{n+1}+v:\ v=\sum_{i=1}^{n}v_iE_{i,n+1}\}\cong\qtwoF^n$ of $P_{n+1}$.

\begin{qtwolemma}[Extension to $\qtwoW(\Pi,\psi^{-1})$ via the Kirillov model]\label{lem:extend}
Let $\qtwoSw(N_n\backslash G_n;\psi^{-1})$ denote the space of locally constant functions
$f:G_n\to\qtwoC$ with $f(ng)=\psi^{-1}(n)f(g)$ for $n\in N_n$, $g\in G_n$, and with support
compact modulo $N_n$. For every such $f$ there exists $W\in\qtwoW(\Pi,\psi^{-1})$ with
$W(\qtwodiag(g,1))=f(g)$ for all $g\in G_n$. In particular, since $W_0$ of \eqref{eq:W0}
lies in $\qtwoSw(N_n\backslash G_n;\psi^{-1})$ \textup{(}it is left $(N_n,\psi^{-1})$-equivariant,
locally constant, and supported on $N_nK_n$, hence compact modulo $N_n$\textup{)}, there is
$W\in\qtwoW(\Pi,\psi^{-1})$ with $W(\qtwodiag(g,1))=W_0(g)$ for all $g\in G_n$.
\end{qtwolemma}

\begin{proof}
This is the Bernstein--Zelevinsky theory of the Kirillov model for $\qtwoGL_{n+1}$, in the
form used by Jacquet--Piatetski-Shapiro--Shalika \cite{JPSS83,Bernstein,Matringe}. We recall
the precise statements and how they combine.

Restriction of Whittaker functions along $\qtwoW(\Pi,\psi^{-1})\to\{W|_{P_{n+1}}\}$ realizes
$\Pi|_{P_{n+1}}$ in a space of functions on $P_{n+1}$; composing with the
homeomorphism $N_n\backslash G_n\xrightarrow{\ \sim\ }(\qtwodiag(G_n,1)\cap N_{n+1}\text{-cosets})$
identifies it with a subspace of functions on $N_n\backslash G_n$ transforming by $\psi^{-1}$
under left $N_n$. This is the \emph{Kirillov model} $\mathcal K(\Pi,\psi^{-1})$ of $\Pi$. The
structural fact we use is:
\begin{itemize}\setlength\itemsep{2pt}
\item[(K1)] \emph{(Bernstein--Zelevinsky derivative/Kirillov theorem; \cite[2.5,~3.5]{Bernstein},
\cite[(2.3)]{JPSS83}.)} The Kirillov space $\mathcal K(\Pi,\psi^{-1})$ contains
$\qtwoSw(N_n\backslash G_n;\psi^{-1})$, and the inclusion is the restriction of $\Pi$ to the
mirabolic $P_{n+1}$. Equivalently: every $f\in\qtwoSw(N_n\backslash G_n;\psi^{-1})$ is the
restriction to $\qtwodiag(G_n,1)$ of some $W\in\qtwoW(\Pi,\psi^{-1})$.
\end{itemize}
Fact (K1) is precisely the assertion of the lemma. We spell out why $W_0$ qualifies:
\eqref{eq:W0} is a finite-level locally constant function (it is bi-invariant under a small
open subgroup of $K_n$ on the right, being defined by integrating $\mathbf 1_{K_n}$ against the
locally constant $\psi$); it is left $(N_n,\psi^{-1})$-equivariant by the change of variables
$x\mapsto nx$ in \eqref{eq:W0}; and it is supported on $N_nK_n$ (shown above), a single
$N_n$-coset family with compact transversal $N_n\backslash N_nK_n\cong (N_n\cap K_n)\backslash K_n$,
hence has compact support modulo $N_n$. Thus $W_0\in\qtwoSw(N_n\backslash G_n;\psi^{-1})$, and (K1)
provides the required $W$.
\end{proof}

\begin{qtwolemma}[Right translation by $u_1$]\label{lem:u1translate}
Let $W\in\qtwoW(\Pi,\psi^{-1})$ be arbitrary, $u_1:=I_{n+1}+E_{n,n+1}\in N_{n+1}$, and
$g\in G_n$ identified with $\qtwodiag(g,1)\in P_{n+1}$. Then for every $h\in G_{n+1}$,
\[
W(\qtwodiag(g,1)\,u_1\,h)=\psi(-g_{nn})\,W(\qtwodiag(g,1)\,h).
\]
In particular, applying this with $h=d_Q^{-1}$ and the extension $W=W_0$ of
Lemma~\ref{lem:extend}, and using $t_Q=u_1 d_Q^{-1}$,
\[
W_0(g\,t_Q)=W_0(\qtwodiag(g,1)\,u_1\,d_Q^{-1})=\psi(-g_{nn})\,W_0(\qtwodiag(g,1)\,d_Q^{-1})
=\psi(-g_{nn})\,W_0(g\,d_Q^{-1}),
\]
the last equality because $\qtwodiag(g,1)\,d_Q^{-1}=\qtwodiag(g\,d_Q^{-1},1)$ lies in
$\qtwodiag(G_n,1)$, where the extension restricts to \eqref{eq:W0}.
\end{qtwolemma}

\begin{proof}
The point is that right multiplication by $u_1$ is converted into \emph{left} multiplication
by an element of $N_{n+1}$, on which $W$ transforms by the generic character $\psi^{-1}$. Set
\[
n_g:=\qtwodiag(g,1)\,u_1\,\qtwodiag(g,1)^{-1}\in G_{n+1}.
\]
Since $u_1=I_{n+1}+E_{n,n+1}$ and conjugation by $\qtwodiag(g,1)$ fixes the $(n+1)$st row
and column index $n+1$ while acting by $g$ on the first $n$ coordinates, a direct computation
gives
\[
n_g=I_{n+1}+\sum_{i=1}^{n}g_{in}\,E_{i,n+1}\in U\subset N_{n+1},
\]
i.e.\ $n_g$ is the upper unipotent matrix whose only nonzero off-diagonal entries are
$(n_g)_{i,n+1}=g_{in}$ for $1\le i\le n$. \textup{(}Explicitly: the $j$th column of
$\qtwodiag(g,1)u_1$ equals that of $\qtwodiag(g,1)$ for $j\le n$, while its $(n+1)$st column is
the $n$th column of $\qtwodiag(g,1)$ plus $e_{n+1}$; multiplying on the right by
$\qtwodiag(g,1)^{-1}$ returns the first $n$ columns to $e_1,\dots,e_n$ and sends the
$(n+1)$st column to $e_{n+1}+\sum_{i=1}^n g_{in}e_i$.\textup{)}

The standard generic character of $N_{n+1}$ is $u\mapsto\psi(\sum_{i=1}^{n}u_{i,i+1})$, depending
only on the superdiagonal entries $u_{i,i+1}$. Among the nonzero entries $(n_g)_{i,n+1}=g_{in}$
of $n_g$, the only one lying on the superdiagonal $\{(i,i+1):1\le i\le n\}$ is the entry
$(n,n+1)$, equal to $g_{nn}$ \textup{(}the entries $(i,n+1)$ with $i<n$ are strictly above the
superdiagonal\textup{)}. Hence the value of the generic character of $\Pi$'s model, namely
$\psi^{-1}$, at $n_g$ is
\[
\psi^{-1}(n_g)=\psi\!\big(-(n_g)_{n,n+1}\big)=\psi(-g_{nn}).
\]
By definition of $n_g$, $\qtwodiag(g,1)\,u_1=n_g\,\qtwodiag(g,1)$, so for any $h$,
\[
W(\qtwodiag(g,1)\,u_1\,h)=W(n_g\,\qtwodiag(g,1)\,h)
=\psi^{-1}(n_g)\,W(\qtwodiag(g,1)\,h)=\psi(-g_{nn})\,W(\qtwodiag(g,1)\,h),
\]
using the left $(N_{n+1},\psi^{-1})$-equivariance of $W\in\qtwoW(\Pi,\psi^{-1})$. The displayed
consequence follows by taking $h=d_Q^{-1}$ and recalling $t_Q=u_1 d_Q^{-1}$
\textup{(}verified directly: $u_1 d_Q^{-1}=(I+E_{n,n+1})\,\qtwodiag(1,\dots,1,Q^{-1},1)$ has
$(n,n+1)$-entry $1$ and equals $d_Q^{-1}u_Q$ since $d_Q^{-1}u_Q$ has the same nonzero
off-diagonal entry $(d_Q^{-1})_{nn}\cdot Q=1$ at position $(n,n+1)$\textup{)}.
\end{proof}

\noindent\textbf{The vector $V$.} Given $\pi$ with conductor $\qtwoqq$, let
$V\in\qtwoW(\pi,\psi)$ be the normalized newvector: the unique $K_1(\qtwoqq)$-invariant
vector with $V(1)=1$, where
$K_1(\qtwoqq)=\{g\in K_n:\text{last row}\equiv(0,\dots,0,1)\bmod\qtwoqq\}$
\cite{JPSS81,Matringe}.

We choose a generator $Q$ of $\qtwoqq^{-1}$, so $|Q|=[\qtwooo:\qtwoqq]=q^{c}\ge1$, where
$\qtwoqq=\qtwopp^{c}$ is the conductor. Recall the
conductor is characterized by the $\varepsilon$-factor:
\begin{equation}\label{eq:eps}
\varepsilon(\tfrac12+s,\pi,\psi)=|Q|^{-s}\,\varepsilon(\tfrac12,\pi,\psi),
\qquad |\varepsilon(\tfrac12,\pi,\psi)|=1 .
\end{equation}

\section*{The Godement--Jacquet engine}

\begin{qtwolemma}[Godement--Jacquet functional equation \cite{GJ72}]\label{lem:GJ}
Let $\pi$ be \emph{any} irreducible admissible representation of $G_n$, let $f$ be a matrix
coefficient of $\pi$ and $\phi\in\qtwoSw(\qtwoMn(\qtwoF))$. The zeta integral
\[
Z(\phi,f,s):=\int_{G_n}\phi(g)f(g)\,|\det g|^{\frac{n-1}{2}+s}\,dg
\]
converges absolutely for $\Re(s)\ge\sigma_0$ \textup{(}some $\sigma_0$ depending only on
$\pi$\textup{)}, and is there a finite $\qtwoC$-linear combination of expressions
$|\det g|^{\frac{n-1}{2}+s}$ integrated against compactly supported data; it continues to a
rational function of $q^{-s}$ such that the normalized ratio
\[
\frac{Z(\phi,f,s)}{L(s,\pi)}\ \text{is a polynomial in }q^{-s}\text{ and }q^{s}
\ \text{\textup{(}in particular holomorphic on all of }\qtwoC\text{\textup{)}.}
\]
It satisfies the functional equation
\[
\gamma(s,\pi,\psi)\,Z(\phi,f,s)=Z(\widehat\phi,f^\vee,1-s),\qquad
\gamma(s,\pi,\psi)=\varepsilon(s,\pi,\psi)\frac{L(1-s,\qtwottilde\pi)}{L(s,\pi)},
\]
where $f^\vee(g)=f(g^{-1})$ and $\widehat\phi$ is the $\psi$-self-dual Fourier
transform on $\qtwoMn(\qtwoF)$. Dually, $Z(\widehat\phi,f^\vee,1-s)$ converges absolutely for
$\Re(s)\le-\sigma_1$ \textup{(}equivalently $\Re(1-s)$ large\textup{)}. The local $L$-factor
$L(s,\pi)=\prod_i(1-a_iq^{-s})^{-1}$ is a finite product, so $1/L(s,\pi)$ is a polynomial in
$q^{-s}$ of degree $\le n$ with constant term $1$; the same holds for $\qtwottilde\pi$. The
identities persist when $f$ is taken to be a Whittaker function of $\pi$. None of these
statements requires $\pi$ to be unitary.
\end{qtwolemma}

Let $\qtwoSe(\qtwoF^\times)$ be the Schwartz--Bruhat functions $\beta$ with
$\beta(xy)=\beta(x)$ whenever $|y|=1$ (i.e.\ $\beta$ depends only on $|x|$),
with Mellin inversion $\beta(y)=\int_{(\sigma)}\tilde\beta(s)|y|^s\,ds$. For such
$\beta$ define $\beta^\sharp$ by
$\beta^\sharp(y)=\int_{(\sigma)}\tilde\beta(s)|y|^{-s}\,ds/\gamma(\tfrac12+s,\pi,\psi)$.

\begin{qtwolemma}\label{lem:beta}
Define $\beta$ by $\tilde\beta(s)=\varepsilon(\tfrac12+s,\pi,\psi)/L(\tfrac12+s,\pi)$.
Then:
\begin{enumerate}\setlength\itemsep{2pt}
\item $\beta$ is supported on $\{|Q|\le|y|\le|Q|q^n\}$ and equals
$\varepsilon(\tfrac12,\pi,\psi)$ on $\{|y|=|Q|\}$;
\item $\beta^\sharp$ has Mellin transform $\widetilde{\beta^\sharp}(s)=1/L(\tfrac12+s,\qtwottilde\pi)$,
is supported on $\{1\le|y|\le q^n\}$, and equals $1$ on $\qtwooo^\times=\{|y|=1\}$.
\end{enumerate}
\end{qtwolemma}

\begin{proof}
By definition of $\gamma$,
$\widetilde{\beta^\sharp}(s)=\tilde\beta(s)/\gamma(\tfrac12+s,\pi,\psi)
=\dfrac{\varepsilon(\tfrac12+s,\pi,\psi)}{L(\tfrac12+s,\pi)}\cdot
\dfrac{L(\tfrac12+s,\pi)}{\varepsilon(\tfrac12+s,\pi,\psi)L(\tfrac12-s,\qtwottilde\pi)}
=\dfrac{1}{L(\tfrac12-s,\qtwottilde\pi)}$;
replacing $s\mapsto-s$ in the inversion contour (legitimate since the integrand is a
polynomial in $q^{\pm s}$) gives the stated form $1/L(\tfrac12+s,\qtwottilde\pi)$.

We now read off the supports. Write $X:=q^{-s}$ and recall
$1/L(\tfrac12+s,\pi)=\prod_{i}(1-a_i\,q^{-\frac12}X)$ is a polynomial in $X$ of degree
$d\le n$ (the degree of the standard $L$-denominator of $\pi$), with constant term $1$.
By \eqref{eq:eps}, $\tilde\beta(s)=\varepsilon(\tfrac12,\pi,\psi)\,|Q|^{-s}\,(1/L(\tfrac12+s,\pi))$,
so $\tilde\beta(s)$ is $\varepsilon(\tfrac12,\pi,\psi)$ times a polynomial in $X$ of degree
$d$, multiplied by $|Q|^{-s}$. Mellin inversion then makes $\beta$ supported precisely on
the shells $|y|\in\{|Q|,|Q|q,\dots,|Q|q^{d}\}\subseteq\{|Q|\le|y|\le|Q|q^{n}\}$, and the
value on the lowest shell $|y|=|Q|$ is the coefficient of $X^{0}$, namely
$\varepsilon(\tfrac12,\pi,\psi)$. This proves (1). Identically, $1/L(\tfrac12+s,\qtwottilde\pi)$
is a polynomial in $X$ of degree $\le n$ with constant term $1$, so $\beta^\sharp$ is
supported on $\{1\le|y|\le q^{n}\}$ and equals $1$ on $|y|=1$. This proves (2).
\end{proof}

\begin{qtwolemma}\label{lem:transfer}
Let $\pi$ be \emph{any} generic irreducible admissible representation of $G_n$ \textup{(}no
unitarity hypothesis\textup{)}. Let $f$ be a Whittaker function of $\pi$ and
$\phi\in\qtwoSw(\qtwoMn(\qtwoF))$. Then both integrals below converge absolutely and
\[
\int_{G_n}\phi(g)f(g)\beta(\det g)|\det g|^{n/2}\,dg
=\int_{G_n}\widehat\phi(g)f^\vee(g)\beta^\sharp(\det g)|\det g|^{n/2}\,dg.
\]
\end{qtwolemma}

\begin{proof}
We work throughout with the non-archimedean Mellin transform: for a function $\beta$ on
$\qtwoF^\times$ depending only on $|y|$, and a contour parameter $\sigma\in\bR$, we write
$\int_{(\sigma)}(\cdots)\,ds:=\frac{\log q}{2\pi i}\int_{\sigma-i\pi/\log q}^{\sigma+i\pi/\log q}(\cdots)\,ds$,
so that $\int_{(\sigma)}\tilde\beta(s)\,|y|^{s}\,ds$ recovers the coefficient extraction
$\beta(y)=\sum_{m}b_m\,\mathbf 1_{|y|=q^{m}}$ from its Mellin transform
$\tilde\beta(s)=\sum_m b_m q^{-ms}$ (a Laurent polynomial in $q^{-s}$), \emph{independently of
$\sigma$}; here $b_m\ne0$ for only finitely many $m$ because $\beta$ is compactly supported on
$\qtwoF^\times$.

\smallskip
\emph{Step 1: absolute convergence of both geometric integrals.} Since $F$ is
non-archimedean, $\phi\in\qtwoSw(\qtwoMn(\qtwoF))$ is locally constant with \emph{compact}
support; let $C\subset\qtwoMn(\qtwoF)$ be a compact set containing $\mathrm{supp}\,\phi$.
By Lemma~\ref{lem:beta}(1), $\beta$ is supported on $\{|\det g|\ge|Q|\}$. Hence the integrand
$\phi(g)f(g)\beta(\det g)|\det g|^{n/2}$ is supported on
$\{g\in G_n:\,g\in C,\ |\det g|\ge|Q|\}$, a \emph{compact} subset of $G_n$: it is
closed and bounded inside $C$, and the lower bound $|\det g|\ge|Q|>0$ keeps it away from the
non-invertible locus $\{\det=0\}$, so it is a compact subset of the open set
$G_n\subset\qtwoMn(\qtwoF)$. On a compact subset of $G_n$ the locally constant function
$\phi(g)f(g)\beta(\det g)|\det g|^{n/2}$ is bounded, so the integral converges absolutely. The
same argument applies verbatim to the right-hand integral: $\widehat\phi$ is again
Schwartz--Bruhat (compact support $C'$), and $\beta^\sharp$ is supported on $\{|\det g|\ge1\}$
by Lemma~\ref{lem:beta}(2), so that integrand is supported on the compact set
$\{g\in C':\,|\det g|\ge1\}\subset G_n$. No unitarity of $\pi$ is used; only the compact support
of $\phi,\widehat\phi$ and the determinant lower bounds on $\beta,\beta^\sharp$.

\smallskip
\emph{Step 2: the left side as a contour integral, $\sigma\gg0$.} Recall from
Lemma~\ref{lem:GJ} that $Z(\phi,f,\tfrac12+s)$ converges absolutely for $\Re(s)\ge\sigma_0$ with
$\sigma_0$ sufficiently large. Fix such a $\sigma>\sigma_0$. Insert the Mellin expansion
$\beta(\det g)=\int_{(\sigma)}\tilde\beta(s)|\det g|^{s}\,ds$, which is the \emph{finite} sum
$\beta(\det g)=\sum_m b_m\mathbf 1_{|\det g|=q^{m}}$, into the left-hand integral. Since this is
a finite sum and $\int_{G_n}|\phi(g)f(g)|\,|\det g|^{\frac n2+\Re(s)}\,dg
=Z(|\phi|,|f|,\tfrac12+\Re(s))<\infty$ for $\Re(s)=\sigma>\sigma_0$, we may interchange the
(finite) Mellin sum with the integral and obtain
\[
\int_{G_n}\phi(g)f(g)\beta(\det g)|\det g|^{n/2}\,dg
=\int_{(\sigma)}\tilde\beta(s)\,Z(\phi,f,\tfrac12+s)\,ds ,\qquad \sigma>\sigma_0,
\tag{$\ast$}
\]
matching $\tfrac{n-1}{2}+(\tfrac12+s)=\tfrac n2+s$. \emph{No central-line contour is used.}

\smallskip
\emph{Step 3: the right side as a contour integral, $\sigma'\ll0$.} By the identical argument
applied to $\widehat\phi$, $f^\vee$, and $\beta^\sharp$ (whose Mellin transform is
$\widetilde{\beta^\sharp}(s)=1/L(\tfrac12+s,\qtwottilde\pi)$, again a Laurent polynomial in
$q^{-s}$), there is $\sigma_1$ such that $Z(\widehat\phi,f^\vee,\tfrac12-s)$ converges
absolutely for $\Re(s)\le-\sigma_1$, and for any $\sigma'<-\sigma_1$,
\[
\int_{G_n}\widehat\phi(g)f^\vee(g)\beta^\sharp(\det g)|\det g|^{n/2}\,dg
=\int_{(\sigma')}\widetilde{\beta^\sharp}(s)\,Z(\widehat\phi,f^\vee,\tfrac12-s)\,ds ,\qquad
\sigma'<-\sigma_1 .
\tag{$\ast\ast$}
\]
Here we used $\beta^\sharp(\det g)=\int_{(\sigma')}\widetilde{\beta^\sharp}(s)|\det g|^{-s}\,ds$;
the sign of the exponent is the one fixed in Lemma~\ref{lem:beta}.

\smallskip
\emph{Step 4: the two integrands are equal and entire.} Define
\[
H(s):=\varepsilon(\tfrac12+s,\pi,\psi)\cdot\frac{Z(\phi,f,\tfrac12+s)}{L(\tfrac12+s,\pi)} .
\]
By Lemma~\ref{lem:GJ}, $Z(\phi,f,\tfrac12+s)/L(\tfrac12+s,\pi)$ is holomorphic in $s$ on all of
$\qtwoC$ (this is part of the Godement--Jacquet theorem and requires \emph{no} unitarity), and
$\varepsilon(\tfrac12+s,\pi,\psi)=|Q|^{-s}\varepsilon(\tfrac12,\pi,\psi)$ is a nowhere-vanishing
holomorphic monomial in $q^{-s}$. Hence $H$ is holomorphic on all of $\qtwoC$; being a finite
$\qtwoC$-linear combination, over the monomials of $\tilde\beta$, of the rational function
$Z(\phi,f,\cdot)/L(\cdot,\pi)$ of $q^{-s}$, the function $H$ is in fact a \emph{Laurent
polynomial in $q^{-s}$}, i.e.\ entire and $\frac{2\pi i}{\log q}$-periodic. Now:
\begin{itemize}\setlength\itemsep{2pt}
\item By definition $\tilde\beta(s)=\varepsilon(\tfrac12+s,\pi,\psi)/L(\tfrac12+s,\pi)$, so the
integrand of $(\ast)$ is
\[
\tilde\beta(s)\,Z(\phi,f,\tfrac12+s)
=\frac{\varepsilon(\tfrac12+s,\pi,\psi)}{L(\tfrac12+s,\pi)}\,Z(\phi,f,\tfrac12+s)=H(s).
\]
\item By Lemma~\ref{lem:GJ}, $Z(\phi,f,\tfrac12+s)=Z(\widehat\phi,f^\vee,\tfrac12-s)/\gamma(\tfrac12+s,\pi,\psi)$
as an identity of meromorphic functions on $\qtwoC$, with
$\gamma(\tfrac12+s,\pi,\psi)=\varepsilon(\tfrac12+s,\pi,\psi)\,L(\tfrac12-s,\qtwottilde\pi)/L(\tfrac12+s,\pi)$.
Using $\widetilde{\beta^\sharp}(s)=\tilde\beta(s)/\gamma(\tfrac12+s,\pi,\psi)$
(Lemma~\ref{lem:beta} and its proof),
\[
\widetilde{\beta^\sharp}(s)\,Z(\widehat\phi,f^\vee,\tfrac12-s)
=\frac{\tilde\beta(s)}{\gamma(\tfrac12+s,\pi,\psi)}\,
\gamma(\tfrac12+s,\pi,\psi)\,Z(\phi,f,\tfrac12+s)
=\tilde\beta(s)\,Z(\phi,f,\tfrac12+s)=H(s).
\]
\end{itemize}
Thus the integrand of $(\ast)$ and the integrand of $(\ast\ast)$ are the \emph{same} entire
function $H(s)$.

\smallskip
\emph{Step 5: contour independence.} Since $H$ is a Laurent polynomial in $q^{-s}$, say
$H(s)=\sum_{k}h_k\,q^{-ks}$ (finite sum), the normalized contour integral satisfies
$\int_{(\sigma)}H(s)\,ds=h_0$ for \emph{every} $\sigma\in\bR$, because
$\int_{(\sigma)}q^{-ks}\,ds=\delta_{k,0}$ independently of $\sigma$ (orthogonality of the
characters $t\mapsto e^{-ikt\log q}$ over one period $t\in[-\pi/\log q,\pi/\log q]$, the
factor $q^{-k\sigma}$ being pulled out and multiplying a vanishing integral when $k\ne0$). In
particular $\int_{(\sigma)}H(s)\,ds=\int_{(\sigma')}H(s)\,ds$. Equivalently: $H$ being entire,
the contour may be shifted from $\Re(s)=\sigma$ to $\Re(s)=\sigma'$ without crossing any pole,
and the two horizontal segments cancel by $\frac{2\pi i}{\log q}$-periodicity of $H$.

\smallskip
Combining Steps 2--5,
\[
\int_{G_n}\phi(g)f(g)\beta(\det g)|\det g|^{n/2}\,dg
\overset{(\ast)}{=}\int_{(\sigma)}H(s)\,ds
=\int_{(\sigma')}H(s)\,ds
\overset{(\ast\ast)}{=}\int_{G_n}\widehat\phi(g)f^\vee(g)\beta^\sharp(\det g)|\det g|^{n/2}\,dg,
\]
which is the assertion, valid for every generic irreducible admissible $\pi$. \qed
\end{proof}

\section*{The local computation}

Set $t_Q:=d_Q^{-1}u_Q=u_1 d_Q^{-1}$. Throughout, $\phi_0:=\mathbf 1_{\qtwoMn(\qtwooo)}$,
whose $\psi$-self-dual Fourier transform is $\widehat{\phi_0}=\phi_0$.

The next lemma is the heart of the computation. Its proof rests on a
\emph{support-disjointness} principle, which we isolate first because it is exactly the
point where a naive ``single-shell'' reading of $\beta$ would be wrong.

\begin{qtwolemma}[Support disjointness]\label{lem:support}
With $\beta$ as in Lemma~\ref{lem:beta} one has the \emph{pointwise} identities, for all
$g\in G_n$,
\begin{align}
\beta(\det g)\,\phi_0(g d_Q^{-1})&=\varepsilon(\tfrac12,\pi,\psi)\,\mathbf 1_{K_n}(g d_Q^{-1}),
\label{eq:supp1}\\
\beta^\sharp(\det g)\,|Q|^{n}\,\phi_0\big(d_Q(g-E_{nn})\big)&=|Q|^{n}\,\mathbf 1_{K_1(\qtwoqq)}(g).
\label{eq:supp2}
\end{align}
\end{qtwolemma}

\begin{proof}
\emph{Proof of \eqref{eq:supp1}.} The factor $\phi_0(g d_Q^{-1})=\mathbf 1_{\qtwoMn(\qtwooo)}(g d_Q^{-1})$
is nonzero only when $g d_Q^{-1}\in\qtwoMn(\qtwooo)$, which forces $\det(g d_Q^{-1})\in\qtwooo$,
i.e.\ $|\det g|\le|\det d_Q|=|Q|$. The factor $\beta(\det g)$ is, by
Lemma~\ref{lem:beta}(1), supported on $|\det g|\ge|Q|$. Hence the product is supported on the
single locus $|\det g|=|Q|$. On that locus:
\begin{itemize}\setlength\itemsep{1pt}
\item if $g d_Q^{-1}\in\qtwoMn(\qtwooo)$ then $|\det(g d_Q^{-1})|=|\det g|/|Q|=1$, so
$g d_Q^{-1}$ is an integral matrix with unit determinant, i.e.\ $g d_Q^{-1}\in K_n$, and
$\beta(\det g)=\varepsilon(\tfrac12,\pi,\psi)$ by Lemma~\ref{lem:beta}(1); thus both sides
equal $\varepsilon(\tfrac12,\pi,\psi)$;
\item if $g d_Q^{-1}\notin\qtwoMn(\qtwooo)$ then $\phi_0(g d_Q^{-1})=0$ and
$\mathbf 1_{K_n}(g d_Q^{-1})=0$ (as $K_n\subset\qtwoMn(\qtwooo)$), so both sides vanish.
\end{itemize}
Off the locus $|\det g|=|Q|$ both sides vanish (the left by support disjointness, the right
because $g d_Q^{-1}\in K_n$ forces $|\det g|=|Q|$). This proves \eqref{eq:supp1}; \emph{crucially,
the higher shells $|\det g|=|Q|q,\dots,|Q|q^{d}$ on which $\beta$ is generically nonzero are
annihilated by $\phi_0(g d_Q^{-1})$, because there $|\det(g d_Q^{-1})|>1$ and $g d_Q^{-1}\notin
\qtwoMn(\qtwooo)$.}

\emph{Proof of \eqref{eq:supp2}.} The factor $\phi_0(d_Q(g-E_{nn}))=\mathbf 1_{\qtwoMn(\qtwooo)}(d_Q(g-E_{nn}))$
is nonzero only when $d_Q(g-E_{nn})\in\qtwoMn(\qtwooo)$. Since $d_Q=\qtwodiag(1,\dots,1,Q)$ scales
the last \emph{row}, this says: $g_{ij}\in\qtwooo$ for $i<n$ (all $j$); $Q\,g_{nj}\in\qtwooo$,
i.e.\ $g_{nj}\in Q^{-1}\qtwooo=\qtwopp^{c}$, for $j<n$; and $Q(g_{nn}-1)\in\qtwooo$, i.e.\
$g_{nn}-1\in\qtwopp^{c}$. In particular every entry of $g$ lies in $\qtwooo$, so $g\in\qtwoMn(\qtwooo)$
and $|\det g|\le1$. The factor $\beta^\sharp(\det g)$ is, by Lemma~\ref{lem:beta}(2), supported on
$|\det g|\ge1$. Hence the product is supported on $|\det g|=1$. There $g\in\qtwoMn(\qtwooo)$ with
unit determinant, i.e.\ $g\in K_n$, with last row $\equiv(0,\dots,0,1)\bmod\qtwopp^{c}=\qtwoqq$; this
is precisely the condition $g\in K_1(\qtwoqq)$, and $\beta^\sharp(\det g)=1$ there by
Lemma~\ref{lem:beta}(2). Conversely every $g\in K_1(\qtwoqq)$ satisfies $d_Q(g-E_{nn})\in\qtwoMn(\qtwooo)$
and $\det g\in\qtwooo^\times$. Off $|\det g|=1$ both sides vanish. This proves \eqref{eq:supp2}.
\end{proof}

\begin{qtwolemma}\label{lem:relations}
Let $\phi(x):=\psi(-x_{nn})\,\phi_0(x d_Q^{-1})$. Then for all $g\in G_n$,
\begin{align}
\int_{N_n}\beta(\det xg)\,\phi(xg)\,\psi(x)\,dx
&=\varepsilon(\tfrac12,\pi,\psi)\,W_0(g\,t_Q),\label{eq:rel1}\\
\beta^\sharp(\det g)\,\widehat\phi(g)&=|Q|^n\,\mathbf 1_{K_1(\qtwoqq)}(g).\label{eq:rel2}
\end{align}
\end{qtwolemma}

\begin{proof}
\emph{Proof of \eqref{eq:rel1}.} Fix $g$ and integrate over $x\in N_n$. Since $x$ is
unipotent, $\det(xg)=\det g$, so $\beta(\det xg)=\beta(\det g)$ is constant in $x$.
Because the last row of $x\in N_n$ is $(0,\dots,0,1)$, we have $(xg)_{nn}=g_{nn}$, hence
\[
\phi(xg)=\psi(-(xg)_{nn})\,\phi_0(xg\,d_Q^{-1})=\psi(-g_{nn})\,\phi_0(xg\,d_Q^{-1}).
\]
Therefore the left-hand side equals
\[
\psi(-g_{nn})\,\beta(\det g)\int_{N_n}\phi_0(x\,g d_Q^{-1})\,\psi(x)\,dx .
\]
We now move $\beta(\det g)$ inside and use \eqref{eq:supp1} pointwise in the variable
$h:=g d_Q^{-1}$. For each $x\in N_n$, applying \eqref{eq:supp1} with $g$ replaced by $xh d_Q=xg$
(note $\det(xg)=\det g$, so $\beta(\det g)=\beta(\det xg)$ and $(xg)d_Q^{-1}=x h$):
\[
\beta(\det g)\,\phi_0(x\,g d_Q^{-1})=\beta(\det xg)\,\phi_0\big((xg)d_Q^{-1}\big)
=\varepsilon(\tfrac12,\pi,\psi)\,\mathbf 1_{K_n}\big((xg)d_Q^{-1}\big)
=\varepsilon(\tfrac12,\pi,\psi)\,\mathbf 1_{K_n}(x\,g d_Q^{-1}).
\]
Hence the left-hand side of \eqref{eq:rel1} equals
\[
\varepsilon(\tfrac12,\pi,\psi)\,\psi(-g_{nn})\int_{N_n}\mathbf 1_{K_n}(x\,g d_Q^{-1})\,\psi(x)\,dx
=\varepsilon(\tfrac12,\pi,\psi)\,\psi(-g_{nn})\,W_0(g d_Q^{-1}),
\]
using \eqref{eq:W0}. Finally, by Lemma~\ref{lem:u1translate} below,
$W_0(g\,t_Q)=W_0(g\,d_Q^{-1}u_Q)=\psi(-g_{nn})\,W_0(g\,d_Q^{-1})$, where $W_0$ now
denotes the $G_{n+1}$-extension of Lemma~\ref{lem:extend} and $g$ is identified with
$\qtwodiag(g,1)$. This gives \eqref{eq:rel1}.

\emph{Proof of \eqref{eq:rel2}.} Compute $\widehat\phi$. With $\phi(x)=\psi(-x_{nn})\phi_0(x d_Q^{-1})$,
the $\psi$-self-dual Fourier transform on $\qtwoMn(\qtwoF)$ gives
\[
\widehat\phi(y)=|Q|^{n}\,\widehat{\phi_0}\big(d_Q(y-E_{nn})\big)=|Q|^{n}\,\phi_0\big(d_Q(y-E_{nn})\big),
\]
where the factor $|Q|^{n}=|\det d_Q|^{-n}$ is the Jacobian of $x\mapsto x d_Q^{-1}$, the substitution
$x\mapsto x d_Q$ turns $\phi_0(x d_Q^{-1})$ into $\phi_0(x)$ and $\psi(-x_{nn})$ into $\psi(-Q x_{nn})$,
and the linear phase $\psi(-Q x_{nn})$ shifts the Fourier variable by $E_{nn}$ after scaling by $d_Q$;
since $\widehat{\phi_0}=\phi_0$ this yields the stated formula. (For $n=1$ this is the elementary identity
$\int_{F}\psi(-x)\mathbf 1_{|x|\le|Q|}\psi(xy)\,dx=|Q|\,\mathbf 1_{|y-1|\le|Q|^{-1}}$.) Therefore
\[
\beta^\sharp(\det g)\,\widehat\phi(g)=\beta^\sharp(\det g)\,|Q|^{n}\,\phi_0\big(d_Q(g-E_{nn})\big)
=|Q|^{n}\,\mathbf 1_{K_1(\qtwoqq)}(g)
\]
by \eqref{eq:supp2}. This proves \eqref{eq:rel2}.
\end{proof}

\section*{A direct unfolding cross-check}

It is instructive (and reassuring) to record that the Rankin--Selberg side of
\eqref{eq:rel1} unfolds to a \emph{single-shell} integral, in agreement with
Lemma~\ref{lem:support}.

\begin{qtwolemma}\label{lem:unfold}
For every $V\in\qtwoW(\pi,\psi)$ and $s\in\qtwoC$,
\[
\int_{N_n\backslash G_n} W_0(g\,t_Q)\,V(g)\,|\det g|^{s-\frac12}\,dg
=\int_{G_n}\psi(-g_{nn})\,\mathbf 1_{K_n}(g d_Q^{-1})\,V(g)\,|\det g|^{s-\frac12}\,dg ,
\]
and the right-hand side is supported on the single shell $|\det g|=|Q|$; it equals
$|Q|^{s-\frac12}\,I$ with $I:=\int_{K_n}\psi(-Q k_{nn})\,V(k d_Q)\,dk$ (Haar on $K_n$ of total
mass $1$).
\end{qtwolemma}

\begin{proof}
Using $W_0(g t_Q)=\psi(-g_{nn})W_0(g d_Q^{-1})$ and \eqref{eq:W0},
\[
W_0(g d_Q^{-1})=\int_{N_n}\mathbf 1_{K_n}(x\,g d_Q^{-1})\,\psi(x)\,dx .
\]
Substitute into the left-hand integral and unfold $N_n\backslash G_n$ against $N_n$: the map
$(x,N_n g)\mapsto xg$ identifies $N_n\times(N_n\backslash G_n)$ with $G_n$ as measure spaces.
Since $(x^{-1}g')_{nn}=g'_{nn}$ for $x\in N_n$ (last row $e_n$), $V(x^{-1}g')=\psi(x^{-1})V(g')
=\psi(-x)V(g')$, and the phase $\psi(x)$ from \eqref{eq:W0} cancels the $\psi(-x)$ from the
$\psi$-equivariance of $V$, one obtains the stated $G_n$-integral. The integrand is supported
where $g d_Q^{-1}\in K_n$, i.e.\ $|\det g|=|Q|$; there $|\det g|^{s-\frac12}=|Q|^{s-\frac12}$ and
$g=k d_Q$ with $k\in K_n$, $g_{nn}=Q k_{nn}$, $dg$ restricting to $dk$, giving $|Q|^{s-\frac12}I$.
\end{proof}

Thus $\int_{N_n\backslash G_n}W_0(g t_Q)V(g)|\det g|^{s-\frac12}\,dg$ is genuinely a single-shell
integral, exactly as Lemma~\ref{lem:support} predicts; multiplying by $\varepsilon(\tfrac12,\pi,\psi)$
and recognizing $\varepsilon\,\mathbf 1_{K_n}(g d_Q^{-1})\psi(-g_{nn})=\int_{N_n}\beta(\det xg)\phi(xg)\psi(x)\,dx$
(this is \eqref{eq:rel1}, re-folded) is what feeds the Godement--Jacquet machinery below. We shall
see that the constant $I$ equals the explicit nonzero scalar of Proposition~\ref{prop:final}.

\section*{Conclusion}

For $W\in\qtwoW(\Pi,\psi^{-1})$, $V\in\qtwoW(\pi,\psi)$, write the Rankin--Selberg
integral
\[
\qtwolRS(s,W,V):=\int_{N_n\backslash G_n}W(\qtwodiag(g,1))\,V(g)\,|\det g|^{s-\frac12}\,dg.
\]

\begin{qtwoproposition}\label{prop:final}
Let $W_0$ be as in \eqref{eq:W0} and $V$ the normalized newvector of $\pi$. Then
for all $s\in\qtwoC$,
\[
\qtwolRS(s,u_Q W_0,\,d_Q V)=c\,|Q|^{-n/2},\qquad
c:=\varepsilon(\tfrac12,\pi,\psi)^{-1}|Q|^n\qtwovol(K_1(\qtwoqq)) .
\]
In particular $\qtwolRS(s,u_Q W_0,d_Q V)=\varepsilon(\tfrac12,\pi,\psi)^{-1}|Q|^{n/2}\qtwovol(K_1(\qtwoqq))$
is a nonzero scalar independent of $s$.
\end{qtwoproposition}

\begin{proof}
A change of variables $g\mapsto g d_Q^{-1}$ gives the homogeneity
$\qtwolRS(s,u_Q W_0,d_Q V)=|Q|^{-(s-\frac12)}\qtwolRS(s,t_Q W_0,V)$, where
$\qtwolRS(s,t_Q W_0,V):=\int_{N_n\backslash G_n}W_0(g t_Q)V(g)|\det g|^{s-\frac12}\,dg$.
By Lemma~\ref{lem:unfold}, $\qtwolRS(s,t_Q W_0,V)=|Q|^{s-\frac12}I$, so it suffices to
evaluate the constant $\varepsilon(\tfrac12,\pi,\psi)\,I$ at any single value of $s$; we work
at $s=\tfrac{n+1}{2}$, i.e.\ with the integrand weight $|\det g|^{n/2}$.

Unfolding via \eqref{eq:rel1} (with $f(g)=V(g)$, $\phi(x)=\psi(-x_{nn})\phi_0(x d_Q^{-1})$),
then applying Lemma~\ref{lem:transfer}, then \eqref{eq:rel2}:
\begin{align*}
\varepsilon(\tfrac12,\pi,\psi)\,\qtwolRS(\tfrac{n+1}{2},t_Q W_0,V)
&=\varepsilon(\tfrac12,\pi,\psi)\!\int_{N_n\backslash G_n}\!W_0(g t_Q)\,V(g)\,|\det g|^{n/2}\,dg\\
&=\int_{N_n\backslash G_n}\!\Big(\int_{N_n}\beta(\det xg)\phi(xg)\psi(x)\,dx\Big)V(g)\,|\det g|^{n/2}\,dg\\
&=\int_{G_n}\phi(g)\,V(g)\,\beta(\det g)\,|\det g|^{n/2}\,dg\\
&=\int_{G_n}\widehat\phi(g)\,V^\vee(g)\,\beta^\sharp(\det g)\,|\det g|^{n/2}\,dg\\
&=|Q|^n\!\int_{K_1(\qtwoqq)}\!V(g^{-1})\,|\det g|^{n/2}\,dg
=|Q|^n\,\qtwovol(K_1(\qtwoqq)).
\end{align*}
The second equality is \eqref{eq:rel1}; the third unfolds the $N_n$-integral into the full
$G_n$-integral exactly as in Lemma~\ref{lem:unfold} (the $\psi$-equivariance of $V$ cancels the
$\psi(x)$ weight), and is the Godement--Jacquet zeta integral $Z(\phi,V,\tfrac12+s)$ at the relevant
point dressed by the Mellin multiplier $\beta$; the fourth is Lemma~\ref{lem:transfer}; the fifth is
\eqref{eq:rel2}, and the last uses that $V$ is right-$K_1(\qtwoqq)$-invariant with $V(1)=1$ and
$|\det g|=1$ on $K_1(\qtwoqq)$, so $V(g^{-1})=1$ there.

Hence $\varepsilon(\tfrac12,\pi,\psi)\,\qtwolRS(\tfrac{n+1}{2},t_Q W_0,V)=|Q|^n\qtwovol(K_1(\qtwoqq))$,
so $\qtwolRS(s,t_Q W_0,V)=c\,|Q|^{s-\frac{n+1}{2}}$ with $c=\varepsilon(\tfrac12,\pi,\psi)^{-1}
|Q|^n\qtwovol(K_1(\qtwoqq))$, and therefore $\qtwolRS(s,u_Q W_0,d_Q V)=|Q|^{-(s-\frac12)}\cdot
c\,|Q|^{s-\frac{n+1}{2}}=c\,|Q|^{-n/2}$. Comparing with Lemma~\ref{lem:unfold} we also identify
$I=\varepsilon(\tfrac12,\pi,\psi)^{-1}|Q|^{n/2}\qtwovol(K_1(\qtwoqq))$.
\end{proof}

Since $\varepsilon(\tfrac12,\pi,\psi)\neq0$, $|Q|>0$, and $\qtwovol(K_1(\qtwoqq))>0$, the scalar
$c\,|Q|^{-n/2}=\varepsilon(\tfrac12,\pi,\psi)^{-1}|Q|^{n/2}\qtwovol(K_1(\qtwoqq))$ is nonzero and
independent of $s$. Hence the Rankin--Selberg integral is finite and nonzero for all $s\in\qtwoC$.
The vector $W:=u_Q W_0$ (built from the single, $\pi$-independent $W_0$) and the $\pi$-newvector
$V$ (rescaled by $d_Q$) exhibit the required test vectors. This proves the Theorem. \qed

\section*{Remarks}

The mechanism is transparent: $u_Q W_0$ is supported on a single $\det$-coset
(Lemma~\ref{lem:unfold}), so the zeta integral collapses to a \emph{constant} in $s$ rather than an
$L$-factor times a unit. What makes the Godement--Jacquet packaging legitimate is
\emph{support disjointness} (Lemma~\ref{lem:support}): although $\beta$ and $\beta^\sharp$ are
genuinely \emph{multi-shell} functions on $\qtwoF^\times$ (with Mellin transforms $1/L(\tfrac12+s,\pi)$
and $1/L(\tfrac12+s,\qtwottilde\pi)$, polynomials of degree up to $n$ in $q^{-s}$), the companion
Schwartz factors $\phi_0(g d_Q^{-1})$ and $\phi_0(d_Q(g-E_{nn}))$ are supported on
$\{|\det g|\le|Q|\}$ and $\{|\det g|\le1\}$ respectively, which meet the supports
$\{|\det g|\ge|Q|\}$ and $\{|\det g|\ge1\}$ of $\beta$ and $\beta^\sharp$ \emph{only on a single shell}.
Thus the higher shells of $\beta$ and $\beta^\sharp$ are annihilated by the Schwartz factors, and the
pointwise identities \eqref{eq:supp1}--\eqref{eq:supp2} hold despite the multi-shell nature of $\beta$.
This is precisely the place where the chosen test pair forces the standard $L$-factors to behave like
units: the full $\beta$ multiplier carries the inverse $L$-factor, but only its leading coefficient
$\varepsilon(\tfrac12,\pi,\psi)$ survives the truncation by $\phi_0$. The $n=1$ case recovers the
classical normalized Gauss sum $I$ and has appeared in subconvexity and moment applications
\cite{Sarnak,MV}; the general $n$ statement, with a single unipotent translate rather than an
average, is the present lemma (joint work in progress with S.~Jana). A version with an average over
many translates appears in \cite{BKL}.

\begin{thebibliography}{9}
\bibitem{Bernstein} J.~N.~Bernstein.
  \emph{$P$-invariant distributions on $\qtwoGL(N)$ and the classification of
  unitary representations of $\qtwoGL(N)$ (non-archimedean case)}. In
  \emph{Lie group representations II}, LNM \textbf{1041}, Springer, 1984,
  50--102.
\bibitem{BKL} A.~R.~Booker, M.~Krishnamurthy, M.~Lee.
  \emph{Test vectors for Rankin--Selberg $L$-functions}.
  J.\ Number Theory \textbf{209} (2020), 37--48.
\bibitem{GJ72} R.~Godement, H.~Jacquet.
  \emph{Zeta functions of simple algebras}. LNM \textbf{260}, Springer, 1972.
\bibitem{JPSS81} H.~Jacquet, I.~I.~Piatetski-Shapiro, J.~Shalika.
  \emph{Conducteur des repr\'esentations du groupe lin\'eaire}.
  Math.\ Ann.\ \textbf{256} (1981), 199--214.
\bibitem{JPSS83} H.~Jacquet, I.~I.~Piatetski-Shapiro, J.~Shalika.
  \emph{Rankin--Selberg convolutions}.
  Amer.\ J.\ Math.\ \textbf{105} (1983), 367--464.
\bibitem{Matringe} N.~Matringe.
  \emph{Essential Whittaker functions for $\qtwoGL(n)$}.
  Doc.\ Math.\ \textbf{18} (2013), 1191--1214.
\bibitem{MV} P.~Michel, A.~Venkatesh.
  \emph{The subconvexity problem for $\qtwoGL_2$}.
  Publ.\ Math.\ IH\'ES \textbf{111} (2010), 171--271.
\bibitem{Sarnak} P.~Sarnak.
  \emph{Fourth moments of Gr\"ossencharakteren zeta functions}.
  Comm.\ Pure Appl.\ Math.\ \textbf{38} (1985), 167--178.
\end{thebibliography}

\endgroup
